% ===== 问题四正文 =====
% 问题四正文；六张示意图的 TikZ 宏在 figures/q4_schematics.tex 中定义。
% 算法依据 Q4_V6_Lite 简化版；实验数字只引用该版本的本地实测记录。
\subsection{问题四模型建立与求解}\label{sec:q4}

问题四和问题三的任务一样，都是从原点出发把 $10$ 到 $16$ 个源全部找到并清除，区别只在源的种类：一部分源是定向的，只朝发射方向两侧各 $90^\circ$ 的范围辐射信号，而且哪些源是定向的、朝哪个方向发射都不知道。这一个改动让第三问里三处依赖“全向接收”的做法都不再成立：一是无信号不再说明离得远，二是覆盖站的布置只考虑了距离，对背对着检测点发射的源不起作用；三是贴着边界朝外发射的源，在场地内部无法接收，只能从场地外侧检测。

第三问中不依赖全向假设的部分（定位区域的半平面交更新、可靠接收区的顶点检验等）在本问仍然成立，只作引用不再重复。

\subsubsection{模型准备}\label{sec:q4-1}

\textbf{（1）定向源的信号模型。}记某个频道的源位于 $p$，接收半径为 $R_f$，它是 $1000$ 到 $1500$ m 之间的一个固定值，但不知道具体是多少。全向源在距离不超过 $R_f$ 的任何位置都能被收到。定向源多一个未知的发射方向，记为单位向量 $u$；它只向 $u$ 两侧各 $90^\circ$ 的范围辐射，也就是只向以“过 $p$ 且垂直于 $u$ 的直线”为边界、包含 $u$ 那一侧的半平面辐射。于是在位置 $x$ 能收到定向源信号的条件是

\begin{equation}\label{eq:q4-1}
\|x-p\|\le R_f\quad\text{且}\quad u\cdot(x-p)\ge0 .
\end{equation}

第二个条件说的就是 $x$ 落在那个半平面里，图\ref{fig:q4-emission}(a) 画出了它的样子。收到信号后返回的示向度、$5$ m 内的近场反馈，以及 $20$ m 内才能完成的光学定位与清除，都和前面几问相同；光学清除不受发射方向影响，只要走到源的 $20$ m 以内就能清除。

\textbf{（2）无信号的两种原因。}对全向源，某处收不到信号只有一个解释：离源超过了 $R_f$。第三问利用这一点排除无信号点周围的保证接收圆盘，并将同频道的有信号点与无信号点配对，得到距离不等式\eqref{eq:q3-3}。对定向源，无信号还可能是因为检测点在发射范围之外，图\ref{fig:q4-emission}(a) 中的两个无信号点就是这两种情况：一个离得近但在源的背面，一个在正面但离得远。既然分不清是哪种原因，一次无信号本身对源的位置就没有任何约束，上述基于无信号的距离约束不能直接沿用。有信号的观测则和以前完全一样：示向度给出 $\pm1^\circ$ 的楔形，收到信号说明距离不超过 $1500$ m，所以第二问建立的定位区域和它的更新方法在本问照常使用。

\textbf{（3）贴着边界朝外发射的源。}把式\eqref{eq:q4-1} 用在场地边界上，就能看到本问最棘手的情形。若源在半径 $1800$ m 的圆周上，$u$ 正好指向圆外，那么发射范围的边界就是圆在该点的切线，场地内部都在切线的另一侧，无法在那里接收到它，见图\ref{fig:q4-emission}(b)。要发现这样的源，检测点必须放到场地外侧去。题目只说干扰源都在目标区域内，没有限制机器狗的活动范围，本文因此允许机器狗在场地边界外不远处停下检测；后面的检测站布局有一环放在半径 $1865$ m 的圆上，就是为了这类源。

\begin{center}
\begin{minipage}{\linewidth}
\centering
\QFourEmissionFigure
\captionof{figure}{定向源只向一侧辐射。(a) 检测点 A 在源的正面且距离够近，能收到信号；检测点 B 虽距离够近但在源的背面，收不到；检测点 C 在正面但超出接收半径，也收不到。无信号无法区分是哪种原因。(b) 源贴在场地边界上朝外发射，场地内部都在源的背面，需在外侧布置检测站。}
\label{fig:q4-emission}
\end{minipage}
\end{center}

\textbf{（4）符号约定 }目标圆域 $\Omega$、频道 $f$、当前位置 $x_t$、定位区域的外包多边形 $\widehat C_f$、包围圆半径 $\rho(\widehat C_f)$ 和圆心 $z(\widehat C_f)$ 都沿用第二、三问的记号；总耗时仍按式\eqref{eq:q3-2} 由移动、换频、检测、光学定位和清除五项累加，评价指标是清除比例和平均定位清除时间 $T/N$。每个已发现频道的记录保存两组位置：收到过信号的位置 $x^+$ 和没有收到信号的位置 $x^-$，它们在后面裁剪定位区域时要用。

\subsubsection{模型的建立}\label{sec:q4-2}

定向源使第三问的发现和定位方法均不再适用。以下沿\textbf{发现——定位——清除——调度——终止}的链条依次建模：先建立对任意发射方向都成立的覆盖条件，再从无信号中恢复定位信息，定位区域满足条件后清除，最后安排多源的执行顺序和途中优化。

\paragraph{保证发现所有源的检测站布局}

\textbf{（1）能发现定向源的覆盖条件。}第三问的覆盖条件只要求每个源位置 $g$ 的 $1000\,\mathrm{m}$ 内存在检测站，即可保证该源被接收。对定向源而言，站在 $1000\,\mathrm{m}$ 内仍不充分——还须落在发射范围内；由于发射方向 $u$ 未知，条件须对所有方向成立：$g$ 附近 $1000\,\mathrm{m}$ 内至少有一个站落在半平面 $u\cdot(x-g)\ge0$ 中。一个充分条件为

\begin{equation}\label{eq:q4-2}
g\in\operatorname{conv}\{s\in P:\ \|s-g\|\le1000\},\qquad\text{对每个 } g\in\Omega ,
\end{equation}

其中 $P$ 是检测站集合，$\operatorname{conv}$ 表示凸包。条件\eqref{eq:q4-2}的充分性可由反证法得到（见图\ref{fig:q4-cover}(b)）：发射范围的边界是一条过 $g$ 的直线；若 $g$ 附近 $1000\,\mathrm{m}$ 内的站全部位于该直线的同一侧，则其凸包也在该侧，与 $g\in\operatorname{conv}$ 矛盾。因此式\eqref{eq:q4-2}保证过 $g$ 的任意直线两侧各至少有一个附近的站，该站距 $g$ 不超过 $1000\,\mathrm{m}$ 且位于发射范围内，由式\eqref{eq:q4-1}可知信号必被接收。全向源仅需附近存在站即可，条件自然满足。

边界朝外辐射的源要求部分检测站设在场外。本文采用原点、内环 $8$ 站（半径 $998\,\mathrm{m}$）和外环 $12$ 站（半径 $1865\,\mathrm{m}$）共 $21$ 个检测站（图\ref{fig:q4-cover}(a)）。外环的设计依据是：相邻两站间隔 $30^\circ$，其连线中点到原点的距离为 $1865\cos15^\circ\approx1801.5\,\mathrm{m}$，恰好使半径 $1800\,\mathrm{m}$ 的场地边界落入凸包。内环取 $998\,\mathrm{m}$ 略小于 $1000\,\mathrm{m}$，使原点附近的点可同时纳入原点和整个内环构成的凸包。条件\eqref{eq:q4-2}的验证采用递归细分方格逐点检验凸包包含关系，共通过 $4228$ 个方格，检验使用保守的浮点比较而非抽样。

扫描方式据此确定：到达每个站时扫描全部未发现的频道。第三问中利用无信号推断跳过的检测在本问不再适用，因为无信号不提供任何几何结论，故对未知频道一律执行检测。$21$ 个站全部扫描完毕后，尚未出现的频道即可判定无源；若已发现 $16$ 个频道（源数上限），则剩余站和频道无须再扫。

\begin{center}
\begin{minipage}{\linewidth}
\centering
\QFourCoverageFigure
\captionof{figure}{$21$ 个检测站的布局与凸包条件。(a) 原点 $1$ 个、内环 $8$ 个、外环 $12$ 个；红色虚线圆是某个源位置 $g$ 的 $1000\,\mathrm{m}$ 范围，绿色区域是范围内各站的凸包。(b) 放大 $g$ 附近：$g$ 在凸包内，无论源朝哪个方向发射，分界线两侧都至少有一个站，发射侧的站一定能收到信号。}
\label{fig:q4-cover}
\end{minipage}
\end{center}

\paragraph{如何从无信号中恢复定位信息}

覆盖条件保证了每个源至少会被一个检测站接收到。接下来的任务是逐步缩小定位区域。有信号时的更新与第三问完全相同；定向源的关键区别在于无信号不再提供距离约束——以下给出两种从无信号中恢复定位信息的方法。

\textbf{（2）有信号时的更新。}频道首次在站 $s$ 收到示向度 $\theta$ 时，定位区域初始化为楔形 $W(s,\theta)$、以 $s$ 为圆心的 $1500\,\mathrm{m}$ 圆盘和场地三者之交，多边形外包方式与第二问一致。此后每收到一次示向度即与当前区域再取交集；若检测设备提示源在 $5\,\mathrm{m}$ 以内，则可直接前往清除。该更新规则对全向源和定向源均适用：收到信号意味着式\eqref{eq:q4-1}的距离和方向条件均已满足，方向条件不再提供额外约束。

\textbf{（3）成对检测点：让无信号也能用。}单个无信号点不提供任何约束，但两个无信号点的联合观测可以推出有用的结论。设上一次收到信号的位置为 $s$，示向方向的单位向量为 $e$，$v$ 是与 $e$ 垂直的单位向量，取

\begin{equation}\label{eq:q4-3}
q^{\pm}=s+t\,e\pm h\,v,\qquad
\frac ht>\tan\delta,\qquad
\Bigl(\frac ht\Bigr)^2+2\,\frac ht\tan\delta<1 ,
\end{equation}

也就是沿示向方向前进 $t$，再向两侧各错开 $h$，见图\ref{fig:q4-pair}。第一个条件保证两点横跨整个 $\pm\delta$ 的楔形；第二个条件保证只要源沿 $e$ 的坐标不小于 $t$，两点到源的距离都比 $s$ 到源的距离小。

由此可以断言：\textbf{若 $q^+$、$q^-$ 都无信号，则源沿 $e$ 的坐标小于 $t$。}反证如下。假设源在 $t$ 之外，由第二个条件，两点都比 $s$ 更靠近源，而 $s$ 已经收到过信号，所以两点都在接收半径之内。另一方面，由第一个条件，源在两点张成的锥形之内，过 $s$ 的任何半平面都会包含 $q^+$、$q^-$ 中的至少一个——因为如果两点都在半平面之外，那么连接 $s$ 和源的线段会穿过 $q^+q^-$，而 $s$ 在半平面内、源也在半平面内（源在发射范围内），中间的穿越点不可能在外面，矛盾。因此至少一个点既在接收范围内，又在发射范围内，不可能同时无信号。这个推断对全向源和定向源都成立，不需要知道 $u$。

两点都无信号时，用半平面 $e\cdot(p-s)\le t$ 裁剪定位区域；若其中一点收到示向度，则按（2）取交集缩小区域，并以该点作为下一轮成对检测的基准。反复执行，每轮沿示向方向截去一段，定位区域逐步缩短。

$t$ 的大小决定每轮截去多少。程序中取当前区域沿 $e$ 的最远距离的 $30\%$；如果区域沿 $e$ 的最近距离已经超过这个值，就把 $t$ 提到最近距离的 $98\%$（但不超过最远距离的 $95\%$），避免检测点落在已知无源的区域。$h$ 取最远距离的 $2.5\%$，并保证 $h/t$ 至少是 $1.05\tan\delta$，此时 $h/t\approx0.08$，式\eqref{eq:q4-3} 的两个条件均满足。实际执行时优先检测距当前位置较近的一点，若收到信号则无须检测另一点。每个源的成对检测至多 $12$ 轮。

\begin{center}
\begin{minipage}{\linewidth}
\centering
\QFourPairFigure
\captionof{figure}{成对检测点的布设与推断逻辑。沿示向方向前进 $t$、向两侧各错开 $h$ 布设一对检测点。若源在 $t$ 之外，两点都比上一个有信号点更靠近源（在接收范围内），且横跨整个楔形（至少一点在发射范围内），因此不可能两点同时无信号——反过来，两点同时无信号就意味着源在 $t$ 之内。}
\label{fig:q4-pair}
\end{minipage}
\end{center}

\textbf{（4）利用无信号点约束发射方向。}（3）利用成对检测点截断定位区域，但没有区分无信号的两种原因：距离太远还是落在源的背面。如果能排除距离太远的可能，就能确定无信号是因为在源的背面，从而限制发射方向的范围，进一步缩小定位区域。思路分三步：先估计接收半径的下界，再判断无信号点是否一定在接收范围内，最后由"在范围内但无信号"推出发射方向的约束。

\emph{估计接收半径的下界。}接收半径至少 $1000$ m，而源到每个有信号点的距离都不超过 $R_f$。由定位区域可以计算 $R_f$ 的一个下界 $R_{\min}$：取 $1000$ m 和”各有信号点到包围圆圆心的距离减去包围圆半径”中的最大者。

\emph{判断检测点在发射范围之外。}若某无信号点 $x^-$ 到定位区域内每一点的距离均小于 $R_{\min}$，则距离超限的可能性被排除，无信号只能是因为 $x^-$ 落在发射范围之外。这同时说明该源是定向的。

\emph{联合约束位置与方向。}确认 $x^-$ 在发射范围之外后，对源的位置 $p$ 和发射方向 $u$ 有：每个有信号点 $x^+$ 满足 $u\cdot(x^+-p)\ge0$，$x^-$ 满足 $u\cdot(x^--p)<0$。两式相减得 $u\cdot(x^+-x^-)>0$，即 $u$ 与向量 $x^+-x^-$ 的夹角小于 $90^\circ$，可能的发射方向被限制在半个圆周内。对每个允许的方向 $u$，$p$ 又被夹在过 $x^-$、$x^+$ 且垂直于 $u$ 的两条平行线之间（图\ref{fig:q4-negative}）。程序将允许的方向等分为 $16$ 个扇区，每个扇区以中心方向执行裁剪并按扇区宽度向外放宽，最后对各扇区的结果取凸包。所得仍为包含真实源的凸多边形，但面积小于原区域。该裁剪仅在 $x^-$ 满足上述条件时执行，否则保持原区域不变。

\begin{center}
\begin{minipage}{\linewidth}
\centering
\QFourNegativeFigure
\captionof{figure}{利用无信号点约束发射方向。无信号点 $x^-$ 到定位区域内每一点的距离都小于接收半径下界 $R_{\min}$，说明 $x^-$ 一定在接收范围内，无信号只能是因为它在源的背面。由此限定发射方向的可能范围，并对每个允许方向用两条垂线将源夹在条带内，条带外的区域被裁掉。}
\label{fig:q4-negative}
\end{minipage}
\end{center}

\paragraph{定位区域满足条件后的清除}

经上述更新，定位区域逐步缩小。当包围圆半径降至清除所需范围时，即进入清除阶段。

\textbf{（5）清除。}定位区域经上述截断和裁剪持续缩小。当区域小到可以从一个位置完成光学清除时，前往该位置执行清除；当区域窄但仍较长时，沿长轴逐点尝试。

具体而言，包围圆半径不超过 $r_c=19.5\,\mathrm{m}$ 时，选取距定位区域各顶点均不超过 $r_c$ 的最近位置执行清除。若检测设备已提示源在 $5\,\mathrm{m}$ 以内，则移动至距该提示点 $14.5\,\mathrm{m}$ 以内即可完成清除。

当定位区域的半宽已小于 $19.5\,\mathrm{m}$ 但长轴仍较长时，可改用逐点光学清除：沿长轴依次摆放半径 $19.5\,\mathrm{m}$ 的圆盘覆盖区域，由近端起依次前往各圆心尝试光学清除，每次失败计 $3\,\mathrm{s}$。源必落在某个圆盘内，故该过程必定成功；仅在所需圆盘不超过 $12$ 个时采用此策略。

\paragraph{多源的执行顺序与途中优化}

以上发现、定位和清除构成处理单个源的完整流程；面对多个源，剩下的问题是执行顺序和途中优化。

\textbf{（6）任务排序与共享测量。}面对多个未清除的源和未访问的检测站，需要决定下一步去哪里。每一步的任务集包含两类节点：未访问的检测站和已发现未清除的源（以包围圆圆心或 $5\,\mathrm{m}$ 提示点表示位置）。任务至多 $21+16=37$ 个。路线规划采用最近邻启发式加 2-opt 改进：分别从距当前位置最近的 $8$ 个任务出发构造最近邻序列，各自经 2-opt 优化后取总路程最短者，仅执行其首任务。每完成一次扫描、成对检测或清除后，任务集和定位区域均已更新，随即重新规划。单次规划仅需数毫秒，可在每步反复调用。

路线中经过已知源附近时，可以顺便补测一次示向度以缩小其定位区域，不增加移动代价。执行共享测量的条件为：该源包围圆半径超过 $19.5\,\mathrm{m}$，圆心距当前位置不超过 $600\,\mathrm{m}$ 加半径，当前位置距该频道此前的检测点不少于 $50\,\mathrm{m}$，且以区域中心和长轴两端附近三个代表位置估算的预期半径缩小量超过 $10\,\mathrm{m}$。估算时根据各代表位置与历史正负观测的几何关系判断该位置是否可能被收到，仅将判断为可接收的位置计入缩小量。清除一个源或完成一轮成对检测后，同样按此规则检查是否应补测。

\textbf{（7）蒙特卡洛 rollout 选择检测动作。}（6）决定了下一步去哪个源，但到达后具体执行什么动作仍有选择余地。默认动作是（3）中的一轮成对检测，但成对检测的 $t$、$h$ 取固定比例，未必是当前状态下的最优选择。

本文采用蒙特卡洛 rollout，在 $13$ 个候选动作（标准及六种不同 $t/h$ 的成对检测，以及六个单点检测）中择优。具体做法是：在定位区域内采样 $24$ 个与历史反馈相容的假想源（位置均匀取样，类型、发射方向和接收半径按均匀分布取样，再以与历史正负反馈的吻合程度加权），对每个假想源分别模拟执行各候选动作直至清除，取 $24$ 次的加权平均耗时作为评分。仅当最优候选的评分比默认动作低 $2\,\mathrm{s}$ 以上时方替换。

\textbf{（8）途中补测。}（6）规划了到下一任务的路线，移动途中通常经过数百米路程。如果路线恰好经过某个已知源附近，在路上停下测一次示向度不增加移动距离，只增加检测时间，却可以缩小该源的定位区域，减少后续处理它所需的轮次（图\ref{fig:q4-transit}）。问题在于：停下来的检测时间能否被后续节省的时间补回？

具体规则是：路段不短于 $200\,\mathrm{m}$ 时，考虑在 $1/4$、$1/2$、$3/4$ 三个位置停下。每个停点上可测的频道须满足：包围圆半径至少 $35\,\mathrm{m}$、圆心到停点不超过 $800\,\mathrm{m}$ 加半径、停点离该频道以前的检测点至少 $60\,\mathrm{m}$，且预计一次测向能把半径缩小 $20\,\mathrm{m}$ 以上。停点测到示向度即与当前区域取交集，无信号则仅记入 $x^-$ 列表，不使用成对检测的推断。

是否停下取决于定位区域的缩小量能否补偿检测时间的增加。这一判断依赖当前状态的诸多细节（区域大小、离源多远、剩余多少站未扫等），难以用简单规则刻画，因此本文采用离线学习的方法：在大量模拟场景的各种状态下，分别记录不停和停下两种策略的总耗时，取二者之差

\begin{equation}\label{eq:q4-4}
Y=T_{\text{不停}}-T_{\text{停}}
\end{equation}

作为停测收益的标签。以停点处的 $57$ 个特征（路段长度、各候选频道的区域大小、预计缩小量、正负观测个数、剩余站数等）拟合 $Y$，得到由 $160$ 棵回归树组成的梯度提升模型。训练样本来自 $2380$ 个模拟场景共 $28359$ 个候选状态，按源数的倒数加权使目标与题目的 $T/N$ 指标一致；特征中不包含真实源的位置、类型或发射方向。

在线运行时，对三个候选停点分别计算特征并预测 $Y$，选取预测值最大且为正的停点执行。为控制风险，每局停测至多 $16$ 次，同一任务不得在连续两段路程中均被安排补测，且在前四个检测站扫描完毕前不启用此规则——实验表明，过早改变路线顺序可能推迟最后几个频道的发现，导致较大的时间损失。

\begin{center}
\begin{minipage}{\linewidth}
\centering
\QFourTransitFigure
\captionof{figure}{移动途中的补测。从当前位置到下一任务点的直线路段上，在 $1/4$、$1/2$、$3/4$ 处考虑停下检测附近的已知源；停在路线上不增加路程，只增加检测时间，是否值得停下由回归树预测的总耗时节省决定。}
\label{fig:q4-transit}
\end{minipage}
\end{center}

\paragraph{何时确认全部源已清除}

上述调度持续进行，直至可以确认所有源已被清除。

\textbf{（9）终止条件。}任务在以下两种情况之一下终止：已真实清除 $16$ 个不同频道的源（源数上限保证无遗漏）；或 $21$ 个站均已扫描完毕且已发现的源全部清除（覆盖条件\eqref{eq:q4-2}保证无遗漏）。仅发现 $16$ 个源但尚未全部清除，或仅站已扫完而已知源未处理完毕，均不满足终止条件。

\subsubsection{模型的求解}\label{sec:q4-3}

模型中没有可以一次解出的静态优化问题，求解就是按上述规则逐步执行：每执行一步，用真实反馈更新定位区域和任务集，再决定下一步。流程见图\ref{fig:q4-flow}，步骤归纳为算法\ref{alg:q4-search}。整套流程只使用四种真实反馈：示向度、$5\,\mathrm{m}$ 以内的近距离提示、无信号和清除成功与否。几何部分保证不漏源、不裁掉真实源、判定能清除时一定能清除；蒙特卡洛 rollout 和回归树评分只在允许的动作里选一个，不改变这些保证。程序用纯 Python 实现，不依赖第三方库，回归树的权重随程序附带，在线运行不需要再训练。

\begin{center}
\begin{minipage}{\linewidth}
\centering
\QFourFlowFigure
\captionof{figure}{问题四的求解流程。每执行一步都重新规划；前四个检测站扫完之前不启用途中停测。}
\label{fig:q4-flow}
\end{minipage}
\end{center}

\par\medskip
\noindent\begin{minipage}{\linewidth}
\small
\refstepcounter{algorithm}\label{alg:q4-search}
\textbf{算法\thealgorithm：混合定向源的滚动搜索与清除}\par\smallskip
\noindent\textbf{输入：}目标圆域 $\Omega$、$20$ 个频道、源数范围 $10$ 到 $16$、$21$ 个检测站、设备动作费用
\begin{enumerate}
\setlength{\itemsep}{2pt}
\setlength{\parskip}{0pt}
\item 在原点扫描全部 $20$ 个频道；收到示向度的频道建立定位区域，无信号的位置记入该频道的 $x^-$ 列表；
\item\label{step:q4-check} 若已真实清除 $16$ 个频道，或 $21$ 站扫完且已发现的源全部清除，结束；
\item 把未访问的站和已发现未清除的源组成任务集，按最近邻加 2-opt 排序，取首任务；
\item 若已扫完至少 $4$ 个站：对去首任务途中的三个候选停点算特征并用回归树预测节省，预测值最大且大于零时前往该停点，检测选定的频道并更新区域，回到步骤\ref{step:q4-check}；
\item 首任务是检测站：前往该站，扫描全部未知频道，对附近满足（6）条件的已知源补测，回到步骤\ref{step:q4-check}；
\item 首任务是源：若包围圆半径不超过 $19.5$ m 或能用不超过 $12$ 个圆盘做光学覆盖，前往清除并按真实反馈登记；否则用蒙特卡洛 rollout 在标准成对检测、其他比例的成对检测和单点检测中选一个执行，收到示向度即与当前区域取交集，两点都无信号则用 $e\cdot(p-s)\le t$ 截断区域，并按（4）检查能否用无信号点再裁剪；回到步骤\ref{step:q4-check}。
\end{enumerate}
\end{minipage}
\par\medskip

\subsubsection{结果分析}\label{sec:q4-4}

\textbf{（1）实验设置与评价指标。}采用本地模拟器生成的 $1140$ 个场景，每场包含 $10$ 至 $16$ 个干扰源。其中 $1000$ 个为随机场景，源在半径 $1770$ m 的圆内均匀分布；另外两组各 $70$ 个，分别为全部源均定向发射，以及全部源靠近场地边界并朝外发射（接收半径取下限 $1000$ m，用于检验难以从场地内部收到信号的情况）。计时按题目规则进行，评价指标沿用第三问的式\eqref{eq:q3-17}。

\textbf{（2）任务完成情况。}本文方案在 $1140$ 个场景中均完成全部清除，共清除 $14806$ 个干扰源。三类场景中每个源的平均耗时分别为 $449.89$、$470.26$ 和 $515.43$ s，结果见表\ref{tab:q4-main}。

\Needspace{40mm}
\begin{center}
\begin{minipage}{\linewidth}
\centering\small
\captionof{table}{本文方案在三类本地场景中的任务完成情况与耗时}
\label{tab:q4-main}
\vspace{4pt}
\begin{tabularx}{\linewidth}{@{}Xrrr@{}}
\toprule
场景类型 & \shortstack{全部清除\\局数} & 清除源数 & \shortstack{平均耗时\\(s/源)} \\
\midrule
随机场景 & 1000/1000 & 12986 & 449.89 \\
全部为定向源 & 70/70 & 910 & 470.26 \\
靠近边界且向外发射 & 70/70 & 910 & 515.43 \\
\bottomrule
\end{tabularx}
\end{minipage}
\end{center}

三类场景均完成全部清除，其中靠近边界且朝外发射一组的平均耗时最高。对此，本文在场外设置检测点，以接收从场地内部难以收到的信号，并结合有信号和无信号的检测结果缩小定位区域。
